\documentclass[a4paper,10pt]{article}
\usepackage[francais]{babel}
\usepackage{url}

\input{letter-icube-uds.sty}

\enCouleur

\begin{document}

\name{Stéphane \textsc{Genaud}}
\fonction{Professeur ENSIIE, mise à disposition Université de Strasbourg}
\telephoneillkirch{+33 (0)3688 54542}
\email{genaud@unistra.fr}

\begin{letter}[Recommandation de Garip \textsc{KUSOGLU} pour une allocation de
  contrat doctoral]%
  {à}%
  {Direction des Ressources Humaines du CNRS, Contrat de doctorant « handicap »}


Garip KUSOGLU  est étudiant  en deuxième année  de Master  Informatique, mention
\textit{Science de  L'Ingénierie Logicielle}  à la  faculté de  Mathématiques et
Informatique de l'université de Strasbourg.  Il effectue sa scolarité en Master,
depuis la première  année, sous un statut d'apprenti. Son  apprentissage se fait
dans le laboratoire ICube sur un travail de recherche dans mon équipe ICPS.  Son
maitre d'apprentissage est Bérenger Bramas, chargé de recherche INRIA. Le rythme
d'alternance mensuel dans sa formation est de trois semaines en apprentissage et
une semaine de cours académiques.
  
  Je connais  Garip depuis sa première  année de Master et  j'ai progressivement
  participé  à son  encadrement sur  son  travail de  recherche. J'ai  également
  cotoyé  récemment Garip  en  tant  qu'enseignant dans un  cours  de Master.

À mes yeux, Garip a les qualités nécessaires pour mener à bien un travail de
doctorat. Il  possède des compétences techniques très solides et une très grande
curiosité intellectuelle. Cette curiosité l'a amené à enrichir son bagage de
connnaissances des langages  de programmation et technologies de l'état de l'art
en informatique. De manière plus générale, Garip est curieux des idées
novatrices dans le domaine.

Dans son travail de recherche, il est également très
opiniatre et absolument dédié à son travail. Je peux témoigner qu'une fois
l'objectif fixé, il travaille sans relâche pour atteindre l'objectif. Il a
également les idées très claires sur son travail, et dans nos discussions de
travail, Garip fait preuve d'une grande honnêteté intellectuelle. Ceci nous permet
d'avoir une relation de travail fructueuse.

Le sujet de thèse sur lequel il porterait sa candidature à pour titre
\textit{Automatic task-based parallelization by source-to-source
  transformations}. C'est un sujet que nous souhaitons approfondir et qui  est la suite de son travail
actuel en apprentissage. Bérenger Bramas et moi même avons eu le temps
d'appréhender les compétences de Garip face à ce travail de recherche et sommes
convaincus qu'il sera à même de réaliser une excellente continuation des travaux
engagés. Nous avons déjà pu publier dans une conférence francophone les bases de
ce travail. 
En conclusion, pour toutes ces raisons, nous soutenons fortement sa candidature.

  Veuillez agrééer, l'expression de mes sentiments respectueux.

\begin{flushright}
\includegraphics[width=.20\textwidth]{signgenaud.jpg} 
\end{flushright}

\end{letter}


\end{document}
